\section{Data}
\label{sec:data}

The training mixture draws from three task families that map
directly onto the deployment-facing capabilities of
Section~\ref{sec:intro}: plain Mandarin / English / multilingual
ASR, hotword-conditioned ASR, and target-speaker ASR (TS-ASR).
Tables~\ref{tab:data_asr}--\ref{tab:data_hotwords} enumerate the
underlying public sources; Section~\ref{subsec:tsasr-synth}
describes the offline TS-ASR synthesiser that produces the
target-speaker training examples.

\subsection{Public ASR Sources}
\label{subsec:asr-sources}

Table~\ref{tab:data_asr} lists the public ASR corpora that feed the
Mandarin, English and code-switching streams. ``Cap.'' is the
per-source utterance cap applied before mixing; ``$r$'' is the repeat
factor (most sources are sampled once; see
Section~\ref{subsec:mixing}). The multilingual long tail is routed
through a separate registry that exposes the same per-source cap and
repeat-factor knobs across languages.

\begin{table}[t]
  \caption{Public ASR sources used for supervised fine-tuning.
    ``Cap.'' is the per-source utterance cap before mixing (``--'' =
    no cap); ``$r$'' is the repeat factor.
    Multilingual sources are routed through the registry described in
    the text and not enumerated here.}
  \label{tab:data_asr}
  \centering
  \small
  \setlength{\tabcolsep}{4pt}
  \begin{tabular}{p{7.8cm}lrr}
    \toprule
    Source & Lang. & Cap. & $r$ \\
    \midrule
    LibriSpeech \cite{panayotov2015librispeech}                                     & en       & --   & 1 \\
    MLS English \cite{pratap2020mls}                                                & en       & 500\,k & 1 \\
    CommonVoice 17 EN (cleaned) \cite{ardila2020commonvoice}                        & en       & 500\,k & 1 \\
    GigaSpeech XL \cite{chen2021gigaspeech}                                         & en       & 1\,M & 1 \\
    National Speech Corpus (Singapore EN)                                           & en       & --   & 1 \\
    Singapore English (community release)                                           & en       & --   & 1 \\
    AISHELL-1 \cite{bu2017aishell}                                                  & zh       & --   & 1 \\
    AISHELL-2 \cite{du2018aishell2}                                                 & zh       & 500\,k & 1 \\
    AISHELL-3 \cite{shi2021aishell3}                                                & zh       & --   & 1 \\
    MagicData Mandarin \cite{magicdata}                                             & zh       & --   & 1 \\
    aidatatang\_200zh                                                                & zh       & --   & 1 \\
    Primewords Chinese 100h                                                          & zh       & --   & 1 \\
    THCHS-30 \cite{wang2015thchs30}                                                 & zh       & --   & 1 \\
    CommonVoice 17 ZH-CN \cite{ardila2020commonvoice}                               & zh       & --   & 1 \\
    WenetSpeech \cite{zhang2022wenetspeech}                                         & zh       & 1\,M & 1 \\
    WenetSpeech4TTS                                                                 & zh       & 500\,k & 1 \\
    Emilia ZH (DNSMOS-filtered) \cite{he2024emilia, reddy2021dnsmos}                & zh       & 500\,k & 1 \\
    KeSpeech \cite{tang2021kespeech}                                                & zh (8 dial.) & 500\,k & 1 \\
    TALCS Mandarin--English code-switch \cite{li2022talcs}                          & zh--en   & --   & 1 \\
    FLEURS (zh, en)                                                                 & multi    & --   & 1 \\
    \emph{Multilingual registry} (10 langs $\times$ \{FLEURS, CommonVoice, YODAS, MLS, Emilia\}) & multi    & per-lang & 1 \\
    \bottomrule
  \end{tabular}
\end{table}

\paragraph{Language tag canonicalisation.}
Source manifests use a mixture of ISO-639-1, ISO-639-3 and locale tags
(``zh'', ``zh-cn'', ``en-us'', ``eng''). The data conversion layer
normalises every per-utterance language tag to a Title-Case
English name (e.g.\ ``zh''\,$\to$\,Chinese, ``en-us''\,$\to$\,English)
across the eighteen languages we currently support
(zh, en, ja, ko, fr, de, es, ru, pt, it, nl, tr, ar, hi, vi, th, id, ms);
unknown values are preserved verbatim.

\subsection{Hotword-Conditioned ASR Sources}
\label{subsec:hotword-sources}

Three of the public ASR datasets carry, on every utterance, a
ground-truth list of rare entities mined offline from the transcript
(Table~\ref{tab:data_hotwords}). At training time we do not re-mix
the audio; the same waveform is paired with the hotword-bearing
supervision to teach the model the prompt-conditioned biasing
behaviour. Two CommonVoice splits with the same annotation are
reserved for evaluation, after applying the per-utterance QA
acceptance flag shipped with each split.

\begin{table}[t]
  \caption{Hotword-bearing supervision splits. The audio is the same
    as the corresponding row in Table~\ref{tab:data_asr}; only the
    supervision metadata is augmented with the ground-truth hotword
    list.}
  \label{tab:data_hotwords}
  \centering
  \small
  \begin{tabular}{llcr}
    \toprule
    Split & Underlying audio & Use & $r$ \\
    \midrule
    aidatatang\_200zh hotwords & aidatatang\_200zh   & train              & 3 \\
    MagicData hotwords         & MagicData Mandarin  & train              & 2 \\
    GigaSpeech hotwords        & GigaSpeech XL       & train (cap 1\,M)   & 1 \\
    \midrule
    CommonVoice 17 ZH hotwords & CommonVoice ZH test & test (QA-filtered) & -- \\
    CommonVoice 17 EN hotwords & CommonVoice EN test & test (QA-filtered) & -- \\
    \bottomrule
  \end{tabular}
\end{table}

\subsection{Augmentation Resources}
\label{subsec:aug-resources}

Two pools of side data are used by the TS-ASR synthesis pipeline
(Section~\ref{subsec:tsasr-synth}): room impulse responses (RIRs) and
additive noise. RIR convolutions are truncated to $1.5$\,s (early
reflections dominate intelligibility) and applied via an FFT-based
overlap-add. Noise is mixed at a configured signal-to-noise ratio
(SNR) and the result is peak-normalised to avoid clipping.

\begin{table}[t]
  \caption{Room impulse response and noise pools used by the TS-ASR
    synthesiser. Within each pool, RIRs are sampled uniformly; noise
    sources are sampled with the listed weights per training example.}
  \label{tab:data_rir_noise}
  \centering
  \small
  \begin{tabular}{p{4.0cm}p{6.0cm}r}
    \toprule
    Resource & Source & Weight \\
    \midrule
    SLR26 sim RIRs       & OpenSLR-26 \cite{ko2017studyrir}                            & uniform \\
    SLR28 sim RIRs       & OpenSLR-28 \cite{ko2017studyrir}                            & uniform \\
    SLR28 real RIRs      & OpenSLR-28 (AIR / RWCP / REVERB-2014) \cite{ko2017studyrir} & uniform \\
    \midrule
    MUSAN noise          & SLR17 \cite{snyder2015musan}                                & 0.30 \\
    MUSAN music          & SLR17 \cite{snyder2015musan}                                & 0.20 \\
    AudioSet road traffic & AudioSet road-traffic subset                               & 0.50 \\
    \bottomrule
  \end{tabular}
\end{table}

\subsection{Subsampling and Mixing}
\label{subsec:mixing}

Mini-batches are formed by uniform multiplexing across all data
sources, after the per-source caps and repeat factors of
Tables~\ref{tab:data_asr}--\ref{tab:data_hotwords} have been applied.
The two knobs implement different objectives:

\begin{itemize}
  \item The per-source utterance \emph{cap} truncates over-large
        corpora -- WenetSpeech and GigaSpeech at $1\,$M utterances;
        Emilia, WenetSpeech4TTS, MLS, AISHELL-2, CommonVoice EN,
        KeSpeech at $500\,$k -- so that no single source dominates
        the mux.
  \item The per-source \emph{repeat factor} $r$ up-samples small
        corpora to a comparable scale: hotword splits are repeated
        $r{=}2$--$3$ times to match the magnitude of the large
        unconditioned ASR streams.
\end{itemize}

The synthesised TS-ASR training set (Section~\ref{subsec:tsasr-synth}) is
added with $r{=}1$ because its $\sim$3\,M positive examples already match
the major ASR streams in size.

\subsection{TS-ASR Synthesis}
\label{subsec:tsasr-synth}

\begin{figure}[tbp]
  \centering
  \includegraphics[width=\linewidth]{figures/tsasr-pipeline.pdf}
  \caption{Offline TS-ASR synthesiser. From a target speaker
    $s_\text{tgt}$ the pipeline draws two different recordings --
    one becomes the $3$--$5\,$s enrollment, the other becomes the
    target waveform. Cross-source interferer sampling adds
    $K\!\in\!\{1,2\}$ competing speakers under
    Eq.~\eqref{eq:mixspeakers}, after which an RIR is applied with
    probability $0.5$ (truncated to $1.5\,$s) and an additive noise
    segment from the MUSAN / AudioSet pool is mixed at
    $\mathrm{SNR}\!\sim\!\mathrm{Uniform}(0,25)\,$dB. Five to ten
    percent of every shard takes the dashed anti-hallucination
    branch: a real enrollment is paired with either a noise-only or
    distractor-only mixture and an empty target transcript, so the
    decoder learns to remain silent when the requested speaker is
    absent.}
  \label{fig:tsasr-pipeline}
\end{figure}

The target-speaker training data is produced by an offline
synthesiser (\Cref{fig:tsasr-pipeline}). Given 16 single-speaker
source manifests (10 Mandarin + 6 English;
Table~\ref{tab:tsasr_sources}), each output sample contains a
target waveform mixed with $K$ interferers, an enrollment utterance
from the same speaker but \emph{a different recording}, and the
target's transcript:

\begin{equation}
  \mathrm{mixed}[n] = s_\text{tgt}[n] + \sum_{k=1}^{K}\,\alpha_k\,s_{\text{int},k}[n - o_k],
  \qquad
  \alpha_k = \frac{\rho_\text{tgt}}{\rho_{\text{int},k}}\cdot 10^{-\mathrm{SNR}_k/20},
  \label{eq:mixspeakers}
\end{equation}

where $\rho$ denotes RMS energy, $o_k$ is a uniformly sampled time
offset that places interferer $k$ inside the target, and the overlap
fraction $|\text{int}_k|/|\text{tgt}|\!\sim\!\mathrm{Uniform}[0.1, 1.0]$
controls how much of the target the interferer covers. The
``clean'' default configuration draws
$K\!\sim\!\mathrm{Cat}(\{1{:}0.8, 2{:}0.2\})$,
$\mathrm{SNR}\!\sim\!\mathcal{N}(5, 7^2)$\,dB clipped to $[-5, 20]$,
and $3$--$5$\,s enrollments. After mixing we apply an RIR with
probability $0.5$ (Table~\ref{tab:data_rir_noise}) and an additive
background noise segment with $\mathrm{SNR}\!\sim\!\mathrm{Uniform}(0, 25)$\,dB.

\begin{table}[t]
  \caption{Single-speaker source manifests fed to the TS-ASR
    synthesiser. The synthesiser drops every speaker with $<\!2$
    utterances and every utterance flagged as failing the per-corpus
    QA check; speakers, languages and source datasets are indexed
    independently to allow cross-source interferer sampling.}
  \label{tab:tsasr_sources}
  \centering
  \small
  \begin{tabular}{lp{6.0cm}p{3.5cm}}
    \toprule
    Lang. & Sources & Coverage \\
    \midrule
    zh (10) & AISHELL-1/2/3, MagicData, KeSpeech, aidatatang\_200zh, Primewords, THCHS-30, CommonVoice ZH, Emilia ZH & full speaker pools, capped per speaker \\
    en (6)  & LibriSpeech (clean-100/360, other-500), MLS English, CommonVoice EN, Emilia EN & same \\
    \bottomrule
  \end{tabular}
\end{table}

\paragraph{Anti-hallucination negatives.}
Five to ten percent of every training shard is replaced by
\emph{negative} samples whose ground-truth transcript is the empty
string. Two sub-types (Table~\ref{tab:tsasr_negatives}) are sampled
with equal weight; both preserve a real enrollment but remove the
target speaker from the mixed channel, training the model to remain
silent when the requested speaker is absent.

\begin{table}[t]
  \caption{Anti-hallucination negative sub-types.}
  \label{tab:tsasr_negatives}
  \centering
  \small
  \setlength{\tabcolsep}{4pt}
  \begin{tabular}{lp{4.5cm}p{6.5cm}}
    \toprule
    Sub-type & Composition & Key parameters \\
    \midrule
    Noise-only       & Real enrollment + 1--2 noise / silence segments & duration $[3,15]$\,s; gain $[-30, 0]$\,dB; silence inclusion probability $0.1$; sources weighted road-traffic $0.6$, MUSAN noise $0.3$, MUSAN music $0.1$ \\
    Distractor-only  & Real enrollment + 1--3 non-target speakers      & SNR $[-5, 20]$\,dB; overlap fraction $[0.1, 1.0]$ \\
    \bottomrule
  \end{tabular}
\end{table}

\paragraph{Output and scale.}
The default budget is $3{,}000{,}000$ training samples, with up to
$200$ samples per speaker. Section~\ref{sec:tsasr} reports
evaluation results on a held-out $10\,$k-example test set generated by
the same pipeline with a held-out random seed to guarantee zero
overlap with training.

\subsection{Hotword Pool Construction at Training Time}
\label{subsec:hotword-pool}

Training-time hotword lists are assembled per sample in three stages:

\begin{enumerate}
  \item \textbf{Real terms.} Each ground-truth hotword is
        independently dropped with probability $0.05$ to simulate
        imperfect retrieval recall.
  \item \textbf{Hard negatives.} An online retriever fetches up to
        $K\!=\!10$ confusable terms per sample using the weighted
        score $0.6\cdot\text{char-overlap} +
        0.4\cdot\text{bigram-Jaccard}$ over a per-corpus hotword
        index, augmented for Mandarin by a syllable-level
        nearest-neighbour search.
  \item \textbf{Random distractors.} The remaining
        $n_\text{distractor}\!\in\![1, 20]$ slots are filled from the
        global hotword pool, with the hard-to-random ratio set so
        that hard negatives saturate first.
\end{enumerate}

The whole hotword prompt line is then kept with probability $0.8$
and dropped otherwise, training the same model to behave gracefully
when no hotword side-channel is available at inference time.

\subsection{Evaluation Datasets and Metrics}
\label{subsec:eval-data}

All evaluations are dispatched through a unified dataset registry,
shared between the offline batch decoder and the online batched
inference server. Table~\ref{tab:eval} lists the test sets
reported in this paper. Per-language WER and CER, mixture error
rate (MER) on code-switching and target-speaker test sets, and
silence recall on negative TS-ASR samples are computed by a single
evaluation harness, so that all numbers in the paper come from the
same metric implementation.

\begin{table}[t]
  \caption{Evaluation sets used in this report. The hotword sweep
    pads each test utterance to $k$ in-prompt hotwords by adding
    random distractors drawn from the same corpus.}
  \label{tab:eval}
  \centering
  \small
  \begin{tabular}{lp{8.5cm}l}
    \toprule
    Task & Test sets & Metric \\
    \midrule
    ASR-zh        & AISHELL-1/2/3, MagicData, KeSpeech, THCHS-30,
                    CommonVoice ZH, WenetSpeech (net / meeting),
                    FLEURS ZH                                                  & CER \\
    ASR-en        & LibriSpeech (test-clean / -other), GigaSpeech, MLS,
                    CommonVoice EN, FLEURS EN                                  & WER \\
    Code-switch   & TALCS test \cite{li2022talcs}                              & MER \\
    Hotword ASR   & CommonVoice ZH / EN hotword splits (sweep
                    $k\!\in\!\{0, 5, 10, 15, 20\}$)                            & CER / WER \\
    TS-ASR        & internal hot+TS test ($n\!=\!6{,}227$, with $5$\,\%
                    silent negatives) + LibriMix 2 / 3 \cite{cosentino2020librimix}
                                                                              & MER + silence \\
    \bottomrule
  \end{tabular}
\end{table}
